% !TeX root = ../papers/nc-p3.tex
\section{純関係のない全分岐の二つの次数配置}
\label{sec:pure-branch-degree-patterns}
\begin{theorem}\label{thm:pure-branch-degree-patterns}
曲面射影と純関係のない標準正則窓で片側分岐があれば、全ての分岐曲線列は
巡回的な記述の起点を適切に取ると、次のいずれかの配置になる。
\begin{enumerate}
\item 順中心は位置 $0,1,2$、逆中心は $0,2,3$ にあり、
順に位置 $0,1,2$ から始まる列の次数は $(1,1),(1,1),(1,3,1)$ である。
\item 四位置の全てに順・逆中心が一つずつあり、四列は全て $(1,1)$ である。
\end{enumerate}
添字移動を同型として割る主張ではない。各場合の動く像は同位置で重複しない。
\end{theorem}
\begin{proof}
定理\ref{thm:branch-pencil-resultant}により二直線列を位置0から始める。
各方向の中心は位置につき一つ以下であり、その占有位置集合 $F,B$ は
$\{0,1\}\subset F,\{2,3\}\subset B$ を満たす。列末尾は $F$ と $B$ の全単射を与える。
初回から終回の動く曲線の間の辺だけを移動として数える。位置 $h$ からの移動は、
$h-1\in F$ または $h+2\in B$ なら高々二つである。前者なら初回の移動以外は
必ずその順中心の反対三次曲線に着き、相異なる移動が一般後点を共有しないからである。
後者なら逆に終回への移動以外は反対三次曲線から出る。
特に12回連続して移動すれば位置1から三回動くため、各列の動く位置数 $m$ は12以下である。

以下は小さい整数の有限列挙である。再現のため全条件を記す。
各 $i\in F$ で $2\le m_i\le12$ を選び、$m_0=2$、終点 $i+m_i+1$ の集合は
重複なしに $B$ とする。位置別移動数に上の条件を課す。
各列の $d_1=d_m=1$、全 $d_j\ge1$、同じ位置の次数の総和は7以下とする。
$m=2$ のとき $i+1\in F,i+2\in B$ が必要である。$m>2$ では $d_2$ は
$i+1\in F$ のとき2、そうでなければ3である。これは $d_2=3-f_1$、次位置の全順中心が
初回直線上にあること、$d_2=1$ なら二直線列となることによる。最後から二番目には逆の条件を課す。
前段の初回・終回以外の移動には、それぞれ到着・出発の次数3を課す。
最後に $d_{-1}=d_0=d_{m+1}=d_{m+2}=0$ とし、$0\le j\le m$ に対する
$d_{j+2}-3d_{j+1}+3d_j-d_{j-1}$ が正なら $j\ge2,\ i+j\in B$、
負なら $j+1<m,\ i+j+1\in F$ とする。これは命題\ref{prop:branch-corridor-degree}による。

$T=(1,1), A=(1,3,1), H=(1,3,3,3,3,3,1)$ と略記すると、
全ループの結果は次の五行である。添付コードは各列の位置ごとの許容次数の直積を取り上記の条件だけで検算する。
\[
\begin{array}{c|c|c}
F&B&\text{開始位置順の列}\\\hline
012&023&(T,T,A)\ \text{または}\ (T,T,H)\\
013&123&(T,A,T)\ \text{または}\ (T,H,T)\\
0123&0123&(T,T,T,T)
\end{array}
\]
$H$ の第三位置の直前・直後は埋め込まれた三次曲線である。直前は初回直線の次の完全三次系、
直後はその位相にある別の初回直線を避ける移動の反対三次曲線である。
実際 $H$ が位置2から始まる場合、第三位置は1で順中心0の初回直線があり、
位置1から始まる場合、第三位置は0で順中心3の初回直線がある。
この直線を中心とする射影を第三曲線に制限すると次数1であるが、
引戻しの固定因子は直前の三次埋込みと一意な前点によって長さ1以下である
（定理\ref{thm:branch-cubic-contact}の最初の射影の議論。初回直線上の逆階数は全点で3である）。
従って第三次数は2以下になり、$H$ を排除する。残る二つの三中心行は巡回表記で一致する。
\end{proof}
